\section{Introduction}\label{sec:intro}

Let $\F$ be an arbitrary field and let $A \in \F^{n \by n}$.  It is classical
that when $A$ is nonsingular, a factorization $A = LU$ with $L$ lower
triangular and $U$ upper triangular exists if and only if every leading
principal submatrix $A[1:k,1:k]$ is nonsingular~\cite{higham_gaussian_2011,golub_matrix_2013}.
The standard remedy when this condition fails is pivoting: for any $A$ there
are permutations $P$ (and $Q$) such that $PA = LU$ or $PAQ = LU$.  Pivoting,
however, changes the problem.  When the row and column indices carry meaning
--- temporal order in autoregressive models, causal order in structural
equation models, or a fixed elimination order imposed by an application ---
the factorization must respect the original indexing, and the correct
question becomes: \emph{for which matrices $A$ does $A = LU$ hold with no
permutations at all?}  We refer to~\cite{darve_lu_2026} for an extended
discussion of this motivation.

Two features make the unpivoted question fundamentally different from the
classical one.  First, $A$ may be singular, or may have singular leading
principal blocks, and yet admit a perfectly good factorization.  Second, the
factors must then be allowed to carry \emph{zeros on their diagonals}: we
require only that $L$ be lower triangular and $U$ upper triangular, not that
they be invertible.  The classical minor criterion says nothing in this
regime.  For example,
\begin{equation}\label{eq:swap-vs-exchange}
  A_1 =
  \begin{pmatrix}
    0 & 1\\
    1 & 0
  \end{pmatrix},
  \qquad
  A_2 =
  \begin{pmatrix}
    0 & 0 & 0\\
    0 & 0 & 1\\
    0 & 1 & 0
  \end{pmatrix}
\end{equation}
are both symmetric with singular leading blocks.  The swap matrix $A_1$ has
no factorization $A_1 = LU$ (from $a_{11} = l_{11}u_{11} = 0$, either the
first row or the first column of $LU$ must vanish).  Yet the larger, more
degenerate matrix $A_2$ factors:
\begin{equation}\label{eq:exchange-factorizations}
  A_2 =
  \underbrace{\begin{pmatrix} 0 & 0 & 0\\ 0 & 1 & 0\\ 1 & 0 & 0 \end{pmatrix}
  \begin{pmatrix} 0 & 1 & 0\\ 0 & 0 & 1\\ 0 & 0 & 0 \end{pmatrix}}_{LU}
  =
  \underbrace{\begin{pmatrix} 0 & 0 & 0\\ 1 & 1 & 0\\ 0 & 1 & 1 \end{pmatrix}
  \begin{pmatrix} -1 & & \\ & 1 & \\ & & -1 \end{pmatrix}
  \begin{pmatrix} 0 & 1 & 0\\ 0 & 1 & 1\\ 0 & 0 & 1 \end{pmatrix}}_{LDL^T} .
\end{equation}
Both factorizations essentially require zero diagonal entries: the leading
zero index of $A_2$ acts as a resource that later off-diagonal structure can
be routed through.  Understanding exactly when such routing is possible is
the subject of this paper.

We answer the question for both the unsymmetric and the symmetric form of
the problem, over an arbitrary field, in a unified way.  Write
$\nl(X) = (\#\text{cols of } X) - \rk(X)$ for the (right) nullity.  Our two
main results are the following (\Cref{thm:lu,thm:ldlt}):
\begin{itemize}
  \item $A \in \F^{n \by n}$ admits a factorization $A = LU$ with $L$ lower
    triangular and $U$ upper triangular, zeros allowed on both diagonals, if
    and only if for every $k = 1, \ldots, n$,
    \[
      \nl(A[1:k,1:k]) \le \nl(A[:,1:k]) + \nl({A[1:k,:]}^T) .
    \]
  \item A symmetric $A \in \F^{n \by n}$ admits a factorization $A = LDL^T$
    with $L$ lower triangular (zeros allowed on the diagonal) and $D$
    diagonal if and only if for every $k = 1, \ldots, n$,
    \[
      \nl(A[1:k,1:k]) \le 2 \, \nl(A[:,1:k]) .
    \]
\end{itemize}
The first criterion is equivalent to the rank inequality of Okunev and
Johnson~\cite{okunev_johnson_2005,johnson_okunev_2025}, restated in the
nullity form introduced in~\cite{darve_lu_2026}.  The second criterion, and
the fact that for symmetric matrices the two criteria coincide --- so that a
symmetric matrix admits an unpivoted $LU$ factorization if and only if it
admits an unpivoted $LDL^T$ factorization --- appear to be new.

\paragraph{The unified method}
Both sufficiency proofs follow the same three-step strategy, which is the
main methodological contribution of this paper.
\begin{enumerate}
  \item \emph{Reduce.}  Every $A \in \F^{n \by n}$ can be written
    $A = M \Pi N$ with $M$ invertible lower triangular, $N$ invertible upper
    triangular, and $\Pi$ a \emph{matching matrix}: a matrix with at most one
    nonzero entry in every row and every column
    (\Cref{lem:reduction-lu}).  Every symmetric $A$ can be written
    $A = M P M^T$ with $M$ invertible lower triangular and $P$ a
    \emph{symmetric matching matrix}, whose nonzero blocks are $1 \by 1$
    diagonal entries and $2 \by 2$ blocks
    $\bigl(\begin{smallmatrix} 0 & \alpha\\ \alpha & \beta
    \end{smallmatrix}\bigr)$ (\Cref{lem:reduction-sym}); over $\mathbb{R}$
    the normal form can be scaled to a \emph{signed matching matrix} with
    entries in $\{0, \pm 1\}$ (\Cref{rem:signed}), recovering the
    unpivoted symmetric elimination of Van Dooren and
    Strang~\cite{vandooren_strang_2014}.  The nonzero pattern of the
    normal form is uniquely determined by $A$ through its rank profile
    (\Cref{prop:canonical}); for invertible $A$ this is the Bruhat
    decomposition, and in general it is the rank profile
    matrix~\cite{dumas_pernet_sultan_2017} and its congruence analogue
    (cf.~\cite{szechtman_2007}).
  \item \emph{Transfer.}  All the nullities appearing in the criteria are
    invariant under $A \mapsto MAN$ with $M, N$ invertible triangular
    (\Cref{lem:transfer}).  Hence $A$ satisfies a criterion if and only if
    its normal form does, and on the normal form the criterion collapses to
    a purely combinatorial \emph{prefix condition} on the matching
    (\Cref{lem:comb-lu,lem:comb-sym}).
  \item \emph{Thread.}  A matching matrix whose matching satisfies the
    prefix condition is factored explicitly.  In the unsymmetric case each
    matched entry is realized by its own rank-one term
    $\ell_t u_t^T$, and the prefix condition is exactly Hall's condition for
    assigning a distinct diagonal slot $t \le \min(i,j)$ to each matched pair
    $(i,j)$; a greedy assignment succeeds (\Cref{thm:lu-matching}).  In the
    symmetric case the rank-one terms $d_t \ell_t \ell_t^T$ are symmetric and
    a single off-diagonal pair cannot be realized in isolation; instead, the
    pairs are threaded into disjoint chains rooted at isolated zero indices,
    and each chain is realized by a telescoping sum of symmetric rank-one
    terms (\Cref{thm:ldlt-matching}).
\end{enumerate}
Necessity of the criteria is a short consequence of Sylvester's rank
inequality; for the symmetric case it follows directly from the unsymmetric
case since $A = LDL^T = L \, (DL^T)$ is itself an $LU$ factorization.

\paragraph{Contributions}
\begin{itemize}
  \item A unified, self-contained, constructive proof method (reduce,
    transfer, thread) that yields necessary and sufficient conditions for
    unpivoted $LU$ and unpivoted $LDL^T$ factorizations of possibly singular
    matrices over an arbitrary field, including fields of characteristic
    $2$.
  \item A necessary and sufficient condition for the existence of an
    unpivoted $LDL^T$ factorization with a truly diagonal $D$ and possibly
    singular $L$, which we believe is new, together with the corollary that
    for symmetric matrices unpivoted $LU$ and $LDL^T$ existence are
    equivalent.
  \item A self-contained account of the matching normal forms and their
    canonicity under triangular equivalence and congruence, connecting the
    existence theory to the generalized Bruhat decomposition and rank
    profile matrix~\cite{elsner_bruhat_1979,malaschonok_2010,dumas_pernet_sultan_2017}
    and to the symmetric canonical forms of Van Dooren--Strang and
    Szechtman~\cite{vandooren_strang_2014,szechtman_2007}, with short
    elimination-style proofs valid over every field.
  \item An explicit $O(n^3)$ algorithm that either produces the
    factorization or certifies its nonexistence, implemented in exact
    arithmetic and validated exhaustively against brute-force enumeration
    over small finite fields (\Cref{sec:validation}).
\end{itemize}

\paragraph{Organization}
\Cref{sec:related} surveys related work.  \Cref{sec:notation} fixes notation.
\Cref{sec:results} states the main theorems and their corollaries.
\Cref{sec:necessity} proves necessity.  \Cref{sec:normalforms} introduces
matching matrices and proves the two reduction lemmas and the canonical-form
proposition.  \Cref{sec:transfer} proves the invariance of the criteria.
\Cref{sec:combinatorics} converts the criteria into prefix conditions on the
matching.  \Cref{sec:explicit} factors the normal forms explicitly.
\Cref{sec:assembly} assembles the proofs of the main theorems and works
through examples.  \Cref{sec:validation} describes the algorithm and its
numerical validation.  The Python implementation is listed in
\Cref{sec:appendix-code}.
